\section*{ADR-003: Tauri Over Electron for Desktop Framework}
\addcontentsline{toc}{section}{ADR-003: Tauri Over Electron for Desktop Framework}

\begin{description}[leftmargin=3cm,labelwidth=2.5cm]
\item[\textbf{Status}] Accepted
\item[\textbf{Date}] February 2024
\item[\textbf{Authors}] Symphony Frontend Team
\end{description}

\subsection*{Summary}

We selected Tauri as Symphony's desktop application framework instead of Electron, despite Electron's larger ecosystem and developer familiarity. This decision affects the entire desktop user experience, application performance, and resource utilization characteristics.

\subsection*{Context}

Symphony requires a desktop framework that supports:

\begin{itemize}
\item Cross-platform deployment (Windows, macOS, Linux)
\item Integration with React-based frontend architecture
\item Low resource consumption for developer productivity tools
\item Native system integration for file operations and OS features
\item Security model appropriate for AI-assisted development workflows
\item Performance characteristics suitable for real-time AI interactions
\end{itemize}

The framework choice significantly impacts Symphony's competitive positioning against existing IDEs, particularly regarding resource usage and startup performance that directly affect developer experience.

\subsection*{Decision}

We chose Tauri 2.x as Symphony's desktop framework, implementing:

\begin{itemize}
\item Rust-based backend with WebView frontend rendering
\item Native API integration through Tauri's command system
\item Security model with granular permission controls
\item Custom protocol handlers for Symphony-specific operations
\item Native system tray and menu integration
\end{itemize}

\subsection*{Rationale}

\subsubsection*{Resource Efficiency}
Tauri applications consume 2-10× less memory than equivalent Electron applications. Symphony's base memory footprint is <150MB compared to 300-800MB for typical Electron IDEs, leaving more resources available for AI model execution.

\subsubsection*{Performance Characteristics}
Tauri's Rust backend provides faster startup times (<1 second vs 3-10 seconds for Electron IDEs) and more responsive native operations. This aligns with Symphony's performance-first philosophy.

\subsubsection*{Security Model}
Tauri's permission-based security model provides granular control over system access, essential for AI-assisted development tools that handle sensitive code and credentials. Each API requires explicit permission declaration.

\subsubsection*{Native Integration}
Tauri provides excellent native system integration without the overhead of bundling Chromium. File system operations, OS notifications, and system APIs perform at native speeds.

\subsubsection*{Rust Ecosystem Alignment}
Tauri's Rust foundation aligns with Symphony's backend architecture, enabling code sharing and consistent development practices across the entire stack.

\subsubsection*{Bundle Size Optimization}
Tauri applications are significantly smaller (10-50MB vs 100-300MB for Electron) due to using system WebView instead of bundling Chromium, improving distribution and installation experience.

\subsubsection*{Trade-offs Accepted}
\begin{itemize}
\item Smaller ecosystem compared to Electron's mature plugin system
\item WebView inconsistencies across different operating systems
\item Limited debugging tools compared to Chromium DevTools
\item Smaller community and fewer learning resources
\item Potential compatibility issues with complex web technologies
\end{itemize}

\subsection*{Alternatives Considered}

\subsubsection*{Electron}
\textbf{Pros}: Mature ecosystem, extensive documentation, large community, consistent Chromium rendering, excellent debugging tools
\textbf{Cons}: High memory usage, slow startup times, large bundle sizes, security concerns, resource competition with AI models
\textbf{Rejected because}: Resource consumption conflicts with AI model requirements and performance expectations

\subsubsection*{Native Development (Qt/GTK)}
\textbf{Pros}: Maximum performance, minimal resource usage, native look and feel, direct system integration
\textbf{Cons}: Platform-specific development, complex UI development, limited web technology integration, higher development costs
\textbf{Rejected because}: Development complexity and web technology integration requirements outweigh performance benefits

\subsubsection*{Progressive Web App (PWA)}
\textbf{Pros}: Cross-platform compatibility, web technology familiarity, automatic updates, no installation required
\textbf{Cons}: Limited system integration, browser security restrictions, offline limitations, inconsistent native features
\textbf{Rejected because}: System integration requirements for file operations and AI model management cannot be met

\subsection*{Consequences}

\subsubsection*{Positive}
\begin{itemize}
\item Achieved target resource usage: <150MB memory footprint vs 300-800MB for Electron alternatives
\item Startup time under 1 second compared to 3-10 seconds for competing IDEs
\item Excellent native system integration with file operations and OS features
\item Strong security model with granular permission controls
\item Smaller application bundles improve distribution and user experience
\item Consistent development experience with Rust backend architecture
\end{itemize}

\subsubsection*{Negative}
\begin{itemize}
\item WebView inconsistencies require platform-specific testing and workarounds
\item Limited ecosystem requires custom implementations for some features
\item Debugging complexity when issues span WebView and Rust backend
\item Smaller community results in fewer third-party solutions and examples
\item Some advanced web features may not be available across all platforms
\end{itemize}

\subsection*{Success Criteria}

\begin{itemize}
\item Memory usage under 200MB for base application (Achieved: <150MB)
\item Startup time under 2 seconds (Achieved: <1 second)
\item Cross-platform feature parity (Achieved: 95\% feature parity)
\item Native system integration functionality (Achieved: Full file system and OS integration)
\item Security model implementation (Achieved: Granular permission system)
\end{itemize}

\subsection*{Risks and Mitigations}

\begin{center}
\begin{tabular}{@{}p{4cm}p{2cm}p{6cm}@{}}
\toprule
\textbf{Risk} & \textbf{Impact} & \textbf{Mitigation} \\
\midrule
WebView inconsistencies & Medium & Complete cross-platform testing and fallback implementations \\
Limited ecosystem & Medium & Custom development and community contribution \\
Debugging complexity & Low & Enhanced logging and development tooling \\
Community support & Low & Active participation in Tauri community and documentation \\
Feature limitations & Medium & Careful feature planning and alternative implementations \\
\bottomrule
\end{tabular}
\end{center}

The Tauri framework choice successfully delivers Symphony's performance and resource efficiency requirements while maintaining cross-platform compatibility and native system integration. This decision contributes significantly to Symphony's competitive advantage in resource-constrained development environments where AI models compete for system resources.